%% Generated by Sphinx.
\def\sphinxdocclass{report}
\documentclass[letterpaper,10pt,english]{sphinxmanual}
\ifdefined\pdfpxdimen
   \let\sphinxpxdimen\pdfpxdimen\else\newdimen\sphinxpxdimen
\fi \sphinxpxdimen=.75bp\relax
\ifdefined\pdfimageresolution
    \pdfimageresolution= \numexpr \dimexpr1in\relax/\sphinxpxdimen\relax
\fi
%% let collapsible pdf bookmarks panel have high depth per default
\PassOptionsToPackage{bookmarksdepth=5}{hyperref}
%% turn off hyperref patch of \index as sphinx.xdy xindy module takes care of
%% suitable \hyperpage mark-up, working around hyperref-xindy incompatibility
\PassOptionsToPackage{hyperindex=false}{hyperref}
%% memoir class requires extra handling
\makeatletter\@ifclassloaded{memoir}
{\ifdefined\memhyperindexfalse\memhyperindexfalse\fi}{}\makeatother

\PassOptionsToPackage{booktabs}{sphinx}
\PassOptionsToPackage{colorrows}{sphinx}

\PassOptionsToPackage{warn}{textcomp}

\catcode`^^^^00a0\active\protected\def^^^^00a0{\leavevmode\nobreak\ }
\usepackage{cmap}
\usepackage{fontspec}
\defaultfontfeatures[\rmfamily,\sffamily,\ttfamily]{}
\usepackage{amsmath,amssymb,amstext}
\usepackage{polyglossia}
\setmainlanguage{english}



\setmainfont{NanumMyeongjo}
\setsansfont{NanumGothic}
\setmonofont{NanumGothicCoding}



\usepackage[Sonny]{fncychap}
\ChNameVar{\Large\normalfont\sffamily}
\ChTitleVar{\Large\normalfont\sffamily}
\usepackage{sphinx}

\fvset{fontsize=\small}
\usepackage{geometry}


% Include hyperref last.
\usepackage{hyperref}
% Fix anchor placement for figures with captions.
\usepackage{hypcap}% it must be loaded after hyperref.
% Set up styles of URL: it should be placed after hyperref.
\urlstyle{same}

\addto\captionsenglish{\renewcommand{\contentsname}{Contents:}}

\usepackage{sphinxmessages}
\setcounter{tocdepth}{1}


\usepackage{fontspec}
\usepackage{kotex}



\title{NTLR800Z\sphinxhyphen{}B1}
\date{2026년 10월 30일}
\release{1.0}
\author{jclee@ntmore.kr}
\newcommand{\sphinxlogo}{\sphinxincludegraphics{logo.png}\par}
\renewcommand{\releasename}{릴리스}
\makeindex
\begin{document}

\pagestyle{empty}
\sphinxmaketitle
\pagestyle{plain}
\sphinxtableofcontents
\pagestyle{normal}
\phantomsection\label{\detokenize{index::doc}}


\sphinxAtStartPar
이 문서는 NTmore에서 개발한 산업용 IoT 게이트웨이 NTLR800Z\sphinxhyphen{}B1의 사용설명서입니다.

\begin{sphinxadmonition}{important}{중요:}
\sphinxAtStartPar
이 문서는 2026년 10월 30일 에 최종적으로 생성 되었습니다.

\sphinxAtStartPar
이 문서는 NTLR800Z\sphinxhyphen{}B1 V00 버전 기준으로 작성 되었습니다.
\end{sphinxadmonition}


\begin{sphinxseealso}{더 보기:}

\sphinxAtStartPar
최신 문서 버전은 v1.0.0 (2026\sphinxhyphen{}09\sphinxhyphen{}04)버전입니다.


\end{sphinxseealso}



\chapter{Contents:}
\label{\detokenize{index:contents}}
\sphinxstepscope


\section{제품 개요}
\label{\detokenize{howto/overview:id1}}\label{\detokenize{howto/overview::doc}}

\subsection{제품 소개 및 사양}
\label{\detokenize{howto/overview:id2}}
\sphinxAtStartPar
\sphinxstylestrong{NTLR800Z\sphinxhyphen{}B1} 은 NTmore에서 개발한 산업용 IoT 게이트웨이 이며 LTE Cat1.bis 를 지원합니다.


\begin{savenotes}\sphinxattablestart
\sphinxthistablewithglobalstyle
\centering
\begin{tabular}[t]{*{3}{\X{1}{3}}}
\sphinxtoprule
\sphinxtableatstartofbodyhook
\sphinxAtStartPar
\sphinxstylestrong{항목}
&\sphinxstartmulticolumn{2}%
\begin{varwidth}[t]{\sphinxcolwidth{2}{3}}
\sphinxAtStartPar
\sphinxstylestrong{설명}
\par
\vskip-\baselineskip\vbox{\hbox{\strut}}\end{varwidth}%
\sphinxstopmulticolumn
\\
\sphinxhline
\sphinxAtStartPar
무선 프로토콜
&\sphinxstartmulticolumn{2}%
\begin{varwidth}[t]{\sphinxcolwidth{2}{3}}
\sphinxAtStartPar
LTE Cat1. bis
\par
\vskip-\baselineskip\vbox{\hbox{\strut}}\end{varwidth}%
\sphinxstopmulticolumn
\\
\sphinxhline
\sphinxAtStartPar
전원 공급 방식
&\sphinxstartmulticolumn{2}%
\begin{varwidth}[t]{\sphinxcolwidth{2}{3}}
\sphinxAtStartPar
DC +5V
\par
\vskip-\baselineskip\vbox{\hbox{\strut}}\end{varwidth}%
\sphinxstopmulticolumn
\\
\sphinxhline
\begin{DUlineblock}{0em}
\item[] 사용주파수대역
\item[] 
\item[] 
\item[] 
\item[] 
\end{DUlineblock}
&
\begin{DUlineblock}{0em}
\item[] 송신
\item[] B1:1920 \textasciitilde{}1980MHz
\item[] B3: 1710 \textasciitilde{} 1785MHz
\item[] B3: 824 \textasciitilde{} 849MHz
\item[] B8: 880 \textasciitilde{} 915Mhz
\end{DUlineblock}
&
\begin{DUlineblock}{0em}
\item[] 수신
\item[] B1: 2110 \textasciitilde{} 2170MHz
\item[] B3: 1805 \textasciitilde{} 1880MHz
\item[] B5: 869 \textasciitilde{} 894MHz
\item[] B8: 925 \textasciitilde{} 960MHz
\end{DUlineblock}
\\
\sphinxhline
\sphinxAtStartPar
안테나 형식
&\sphinxstartmulticolumn{2}%
\begin{varwidth}[t]{\sphinxcolwidth{2}{3}}
\sphinxAtStartPar
내장 안테나
\par
\vskip-\baselineskip\vbox{\hbox{\strut}}\end{varwidth}%
\sphinxstopmulticolumn
\\
\sphinxhline
\sphinxAtStartPar
인터페이스
&\sphinxstartmulticolumn{2}%
\begin{varwidth}[t]{\sphinxcolwidth{2}{3}}
\sphinxAtStartPar
RJ45 Connector x2, 6pin Connector, USIM slot
\par
\vskip-\baselineskip\vbox{\hbox{\strut}}\end{varwidth}%
\sphinxstopmulticolumn
\\
\sphinxbottomrule
\end{tabular}
\sphinxtableafterendhook\par
\sphinxattableend\end{savenotes}

\begin{sphinxadmonition}{note}{참고:}
\sphinxAtStartPar
제품 메뉴얼은 제품 출시 후 펌웨어 업그레이드에 따라 내용이 달라질 수 있으니, 초기화면에 \sphinxstylestrong{firmware version} 을 확인해 주십시요.
\end{sphinxadmonition}


\subsection{LED 인디케이터 정의}
\label{\detokenize{howto/overview:led}}

\begin{savenotes}\sphinxattablestart
\sphinxthistablewithglobalstyle
\centering
\begin{tabular}[t]{*{3}{\X{1}{3}}}
\sphinxtoprule
\sphinxtableatstartofbodyhook
\sphinxAtStartPar
\sphinxstylestrong{종류}
&
\sphinxAtStartPar
\sphinxstylestrong{설명}
&
\sphinxAtStartPar
\sphinxstylestrong{설명}
\\
\sphinxhline
\sphinxAtStartPar
\sphinxstylestrong{POWER}
&
\sphinxAtStartPar
전원이 켜짐
&
\sphinxAtStartPar
전원(RED) 켜짐
\\
\sphinxhline
\begin{DUlineblock}{0em}
\item[] \sphinxstylestrong{LAN1,LAN2}
\item[] 
\end{DUlineblock}
&
\begin{DUlineblock}{0em}
\item[] LAN 연결
\item[] LAN 연결 및 소프트웨어 준비
\end{DUlineblock}
&
\begin{DUlineblock}{0em}
\item[] LAN(GREEN) 깜빡임
\item[] LAN(ORANGE) 깜빡임
\end{DUlineblock}
\\
\sphinxhline
\begin{DUlineblock}{0em}
\item[] \sphinxstylestrong{LTE 신호}
\item[] 
\item[] 
\item[] 
\end{DUlineblock}
&
\begin{DUlineblock}{0em}
\item[] 강전계 및 연결 완료(RSRP ≥ \sphinxhyphen{}100dBm)
\item[] 약전계 및 연결 완료(RSRP ≤ \sphinxhyphen{}100dBm)
\item[] 망 연결 실패
\item[] USIM 오류
\end{DUlineblock}
&
\begin{DUlineblock}{0em}
\item[] LTE(GREEN) 켜짐
\item[] LTE(ORANGE) 켜짐
\item[] LTE(ORANGE) 깜빡임
\item[] LTE(RED) 깜빡임
\end{DUlineblock}
\\
\sphinxbottomrule
\end{tabular}
\sphinxtableafterendhook\par
\sphinxattableend\end{savenotes}

\sphinxstepscope


\section{관리자 UI 접속}
\label{\detokenize{howto/administrator:ui}}\label{\detokenize{howto/administrator::doc}}

\subsection{UI 접속}
\label{\detokenize{howto/administrator:id1}}\begin{enumerate}
\sphinxsetlistlabels{\arabic}{enumi}{enumii}{}{.}%
\item {} 
\sphinxAtStartPar
랜 케이블을 이용하여 모뎀과 PC를 연결합니다.

\item {} 
\sphinxAtStartPar
아래의 URL을 웹 브라우저를 사용하여 접속합니다.

\item {} 
\sphinxAtStartPar
아래 그림처럼 로그인 페이지가 나오면 초기 설정되어 있는 사용자 이름과 비밀번호를 입력하세요.

\end{enumerate}
\begin{itemize}
\item {} 
\sphinxAtStartPar
\sphinxurl{http://192.168.15.1}

\end{itemize}

\noindent\sphinxincludegraphics[width=0.500\linewidth]{{login-00}.png}


\subsection{기본 로그인 및 비밀번호 변경}
\label{\detokenize{howto/administrator:id2}}
\noindent\sphinxincludegraphics[width=0.800\linewidth]{{password-00}.png}

\sphinxstepscope


\section{통신사 망 연동 설정}
\label{\detokenize{howto/connection:id1}}\label{\detokenize{howto/connection::doc}}

\subsection{SIM 카드 장착 및 상태 확인}
\label{\detokenize{howto/connection:sim}}
\noindent\sphinxincludegraphics[width=1.000\linewidth]{{modem-00}.png}

\sphinxAtStartPar
\sphinxstylestrong{모뎀\sphinxhyphen{}>모뎀 정보} 화면에서 SIM 상태를 확인 할 수 있습니다.

\sphinxAtStartPar
\sphinxstylestrong{SIM 상태} 가 올바르게 \sphinxstyleemphasis{Registered} 상태인지 확인하십시요.

\sphinxAtStartPar
SIM에 대한 자세한 정보는 해당 메뉴에서 USIM 모양의 아이콘을 클릭 하면 표시됩니다.

\noindent\sphinxincludegraphics[width=1.000\linewidth]{{modem-03}.png}


\subsection{APN 프로필 설정}
\label{\detokenize{howto/connection:apn}}
\noindent\sphinxincludegraphics[width=1.000\linewidth]{{modem-00}.png}

\sphinxAtStartPar
\sphinxstylestrong{모뎀\sphinxhyphen{}>모뎀 정보\sphinxhyphen{}>설정} 화면에서 사용자가 직접 APN 을 설정 할 수 있습니다.

\sphinxAtStartPar
기본값은 NULL (자동 할당) 이며, 이에대한 표시는 \sphinxstylestrong{모뎀\sphinxhyphen{}>모뎀 정보} 화면에 표시 됩니다.

\noindent\sphinxincludegraphics[width=1.000\linewidth]{{modem-00}.png}

\sphinxAtStartPar
\sphinxstylestrong{APN} LTE망에서 받은 APN 값 입니다. NULL 로 설정할 경우 \sphinxstylestrong{|} 가 표시되며, 사용자가 설정한 값 NULL이 표시 됩니다.


\subsection{신호 세기 모니터링}
\label{\detokenize{howto/connection:id2}}
\noindent\sphinxincludegraphics[width=1.000\linewidth]{{modem-04}.png}

\sphinxAtStartPar
신호 세기 모니터링은 \sphinxstylestrong{모뎀\sphinxhyphen{}>모뎀 정보} 하단에 연관된 값들을 보여줍니다.

\sphinxstepscope


\section{장치 관리 및 식별}
\label{\detokenize{howto/device_manage:id1}}\label{\detokenize{howto/device_manage::doc}}

\subsection{단말 정보 확인}
\label{\detokenize{howto/device_manage:id2}}
\sphinxAtStartPar
단말의 정보는 \sphinxstylestrong{상태\sphinxhyphen{}>개요} 에 표시됩니다.

\noindent\sphinxincludegraphics[width=1.000\linewidth]{{overview-01}.png}

\sphinxAtStartPar
NTLR800Z\sphinxhyphen{}B1모델은 디바이스의 고유 번호인 \sphinxstylestrong{디바이스 일련번호} 가 기본이 되고 있습니다.


\subsection{펌웨어 업데이트 방법}
\label{\detokenize{howto/device_manage:id3}}
\sphinxAtStartPar
펌웨어 업데이트는 자사의 FOTA 서버를 통해서만 진행이 가능합니다.

\sphinxAtStartPar
기본 화면은 아래와 같습니다.

\noindent\sphinxincludegraphics[width=1.000\linewidth]{{howto-00}.png}
\begin{enumerate}
\sphinxsetlistlabels{\arabic}{enumi}{enumii}{}{.}%
\item {} 
\sphinxAtStartPar
FOTA 연동상태를 확인하여야 하며 \sphinxstylestrong{비활성화} 인 경우는 \sphinxstylestrong{시스템\sphinxhyphen{}>시스템 업데이트\sphinxhyphen{}>설정} 에서 \sphinxstylestrong{FOTA 기능 활성화} 를 선택하십시요.

\item {} 
\sphinxAtStartPar
\sphinxstylestrong{FOTA 기능 활성화} 를 선택한 후에는 \sphinxstylestrong{연동 확인**과 **즉시 확인} 버튼이 활성화 됩니다.

\item {} 
\sphinxAtStartPar
\sphinxstylestrong{연동 확인} 혹은 {\color{red}\bfseries{}**}즉시 확인* 버튼을 누르면 아래처럼 활성화가 됩니다.

\end{enumerate}

\noindent\sphinxincludegraphics[width=1.000\linewidth]{{howto-01}.png}
\begin{enumerate}
\sphinxsetlistlabels{\arabic}{enumi}{enumii}{}{.}%
\setcounter{enumi}{3}
\item {} 
\sphinxAtStartPar
만약에 자동으로 업데이트를 원하면 \sphinxstylestrong{시스템\sphinxhyphen{}>시스템 업데이트\sphinxhyphen{}>설정} 에서 \sphinxstylestrong{자동 업데이트 설정} 을 활성화 하시면 됩니다.

\item {} 
\sphinxAtStartPar
서버에서 업데이트 내용이 있으면, 업데이트가 진행 됩니다.

\end{enumerate}

\begin{sphinxadmonition}{danger}{위험:}
\sphinxAtStartPar
업데이트가 진행중에는 전원을 끄거나 USIM을 제거하는 동작을 해서는 절대 안됩니다.
\end{sphinxadmonition}


\subsection{공장 초기화}
\label{\detokenize{howto/device_manage:id6}}
\sphinxAtStartPar
공장 초기화는 두가지 방법을 사용할수 있습니다.

\sphinxAtStartPar
\sphinxstylestrong{관리 UI 접속이 가능할때}

\noindent\sphinxincludegraphics[width=0.800\linewidth]{{howto-02}.png}
\begin{enumerate}
\sphinxsetlistlabels{\arabic}{enumi}{enumii}{}{.}%
\item {} 
\sphinxAtStartPar
\sphinxstylestrong{시스템\sphinxhyphen{}>Backup / Restore Config} 에서 \sphinxstylestrong{복원} 기본값으로 초기화 \sphinxcode{\sphinxupquote{초기화 수행}} 을 실행 합니다.

\item {} 
\sphinxAtStartPar
초기 설정 화면 (\sphinxurl{http://192.168.15.1}) 으로 접속하여 로그인 합니다.

\end{enumerate}

\sphinxAtStartPar
\sphinxstylestrong{관리 UI 접속이 가능하지 않은 경우}
\begin{enumerate}
\sphinxsetlistlabels{\arabic}{enumi}{enumii}{}{.}%
\item {} 
\sphinxAtStartPar
전원을 연결하고 디바이스를 뒤집어 봅니다.

\item {} 
\sphinxAtStartPar
reset 이라고 써있는 부분을 약 5초간 누르면, LTE LED가 점등을 시작하며, 빠르게 점등이 가속되면 누르고 있는것을 떼면 됩니다.

\item {} 
\sphinxAtStartPar
약 5분정도 지난후에 초기 설정 화면(\sphinxurl{http://192.168.15.1}) 으로 접속하여 로그인 합니다.

\end{enumerate}

\begin{sphinxadmonition}{danger}{위험:}
\sphinxAtStartPar
초기화가 진행중에 일부러 전원을 자주 끄는 행동을 해서는 안됩니다.
\end{sphinxadmonition}

\sphinxstepscope


\section{로그인}
\label{\detokenize{main_menu/login:id1}}\label{\detokenize{main_menu/login::doc}}

\bigskip\hrule\bigskip

\begin{enumerate}
\sphinxsetlistlabels{\arabic}{enumi}{enumii}{}{.}%
\item {} 
\sphinxAtStartPar
랜 케이블을 이용하여 모뎀과 PC를 연결합니다.

\item {} 
\sphinxAtStartPar
아래의 URL을 웹 브라우저를 사용하여 접속합니다.

\item {} 
\sphinxAtStartPar
아래 그림처럼 로그인 페이지가 나오면 초기 설정되어 있는 사용자 이름과 비밀번호를 입력하세요.

\end{enumerate}
\begin{itemize}
\item {} 
\sphinxAtStartPar
\sphinxurl{http://192.168.15.1}

\end{itemize}

\noindent\sphinxincludegraphics[width=0.500\linewidth]{{login-00}.png}

\begin{sphinxadmonition}{note}{참고:}
\sphinxAtStartPar
패스워드는 디바이스 뒷면에 적혀 있습니다.
\end{sphinxadmonition}

\sphinxstepscope


\section{상태}
\label{\detokenize{main_menu/status:id1}}\label{\detokenize{main_menu/status::doc}}

\subsection{개요}
\label{\detokenize{main_menu/status:id2}}
\sphinxAtStartPar
로그인이 성공하면 제일 처음으로 보여지는 페이지 입니다. 대략적인 디바이스에 대한 정보를 확인 할 수 있습니다.


\subsubsection{Modem status}
\label{\detokenize{main_menu/status:modem-status}}
\noindent\sphinxincludegraphics[width=1.000\linewidth]{{overview-00}.png}

\sphinxAtStartPar
모뎀의 정보를 보여 줍니다. 자세한 항목은 \sphinxstylestrong{모뎀} 항목에서 볼수 있습니다.


\subsubsection{시스템}
\label{\detokenize{main_menu/status:id3}}
\noindent\sphinxincludegraphics[width=1.000\linewidth]{{overview-01}.png}

\sphinxAtStartPar
라우터에 대한 정보를 보여줍니다.


\subsubsection{메모리}
\label{\detokenize{main_menu/status:id4}}
\noindent\sphinxincludegraphics[width=1.000\linewidth]{{overview-02}.png}

\sphinxAtStartPar
메모리 사용량을 보여줍니다.


\subsubsection{네트워크}
\label{\detokenize{main_menu/status:id5}}
\noindent\sphinxincludegraphics[width=1.000\linewidth]{{overview-03}.png}

\sphinxAtStartPar
물리적으로 연결된 모뎀의 정보를 보여줍니다.


\subsubsection{DHCP 할당}
\label{\detokenize{main_menu/status:dhcp}}
\noindent\sphinxincludegraphics[width=1.000\linewidth]{{overview-04}.png}

\sphinxAtStartPar
DHCP 할당에 대한 정보를 보여줍니다. \sphinxguilabel{IP 예약} 을 선택시에 호스트에 대해서 내부 고정 IP 를 설정 합니다.


\subsection{실시간 그래프}
\label{\detokenize{main_menu/status:id6}}
\sphinxAtStartPar
실시간으로 디바이스의 현황을 그래프로 보여줍니다.


\subsubsection{부하}
\label{\detokenize{main_menu/status:id7}}
\noindent\sphinxincludegraphics[width=1.000\linewidth]{{graph-00}.png}

\sphinxAtStartPar
부하 평균은 Linux에서 시스템 자원 현황을 모니터링하기 위해 사용하는 지표입니다.


\subsubsection{대역폭}
\label{\detokenize{main_menu/status:id8}}
\noindent\sphinxincludegraphics[width=1.000\linewidth]{{graph-01}.png}

\sphinxAtStartPar
사용 가능한 모든 물리적 인터페이스의 대역폭 사용 현황을 표시합니다.


\subsubsection{연결}
\label{\detokenize{main_menu/status:id9}}
\noindent\sphinxincludegraphics[width=1.000\linewidth]{{graph-02}.png}

\sphinxAtStartPar
이 장치를 통해 활성화된 연결 현황을 표시합니다. \sphinxguilabel{DNS 조회 사용} 버튼을 클릭하면 목적지까지의 전송량을 표시해 줍니다.

\sphinxstepscope


\section{모뎀}
\label{\detokenize{main_menu/modem:id1}}\label{\detokenize{main_menu/modem::doc}}

\subsection{모뎀 정보}
\label{\detokenize{main_menu/modem:id2}}

\subsubsection{기본 정보}
\label{\detokenize{main_menu/modem:id3}}
\noindent\sphinxincludegraphics[width=0.800\linewidth]{{modem-00}.png}

\sphinxAtStartPar
\sphinxstylestrong{신호 세기} : 신호 세기를 안테나와 \%와 함께 표시 됩니다.

\sphinxAtStartPar
\sphinxstylestrong{통신사} : 연결된 통신사 이름을 표시 합니다.

\sphinxAtStartPar
\sphinxstylestrong{SIM 상태} : SIM 카드의 상태를 표시 합니다.

\sphinxAtStartPar
USIM의 그림을 클릭하면 SIM 카드에 대한 정보를 표시합니다.

\begin{figure}[htbp]
\centering

\noindent\sphinxincludegraphics[width=0.800\linewidth]{{modem-03}.png}
\end{figure}

\sphinxAtStartPar
\sphinxstylestrong{연결 통계} : 연결된 네트워크의 통계 정보를 표시 합니다.

\sphinxAtStartPar
\sphinxstylestrong{통신 기술} : 사용 중인 통신 기술을 표시 합니다.

\sphinxAtStartPar
\sphinxstylestrong{APN} : 사용 중인 APN을 표시 합니다.


\subsubsection{모뎀 식별 정보}
\label{\detokenize{main_menu/modem:id4}}
\noindent\sphinxincludegraphics[width=0.800\linewidth]{{modem-02}.png}

\sphinxAtStartPar
\sphinxstylestrong{모델명} : 사용중인 모뎀의 모델명을 표시 합니다.

\sphinxAtStartPar
\sphinxstylestrong{펌웨어 버전} :   사용중인 모뎀의 펌웨어 버전을 표시 합니다.

\sphinxAtStartPar
\sphinxstylestrong{IP 주소} :   IP를 표시 합니다.

\begin{sphinxadmonition}{important}{중요:}
\sphinxAtStartPar
모뎀은 일반적으로 CGNAT 구간의 IP를 할당받습니다.
\end{sphinxadmonition}


\subsubsection{신호 품질 상세}
\label{\detokenize{main_menu/modem:id5}}
\noindent\sphinxincludegraphics[width=0.800\linewidth]{{modem-04}.png}

\sphinxAtStartPar
신호 품질에 대한 상세 정보를 표시 합니다.


\subsubsection{설정}
\label{\detokenize{main_menu/modem:id6}}
\noindent\sphinxincludegraphics[width=0.800\linewidth]{{modem-01}.png}

\sphinxAtStartPar
\sphinxstylestrong{현재 할당된 네트워크 APN} : 현재 LTE 망 접속을 통해 할당받은 APN 정보를 표시 합니다.

\sphinxAtStartPar
\sphinxstylestrong{사용자 지정 APN} : 사용자가 직접 APN을 입력하여 연결할 수 있습니다. (\sphinxstylestrong{NULL} 입력시 망에서 부여받습니다.)

\sphinxAtStartPar
\sphinxstylestrong{LTE BAND 선택} : 디바이스가 사용가능한 BAND를 선택할 수 있습니다.

\begin{sphinxadmonition}{important}{중요:}
\sphinxAtStartPar
통신사에 따라 LTE BAND 선택이 제한될 수 있습니다.
\end{sphinxadmonition}


\subsubsection{DFota(Delta Firmware Over\sphinxhyphen{}The\sphinxhyphen{}Air)}
\label{\detokenize{main_menu/modem:dfota-delta-firmware-over-the-air}}
\sphinxAtStartPar
모뎀(Quectel)의 FOTA 펌웨어 업데이트 기능을 설정합니다.

\noindent\sphinxincludegraphics[width=0.800\linewidth]{{dfota-00}.png}

\sphinxAtStartPar
\sphinxstylestrong{DFOTA 기능 활성화} : DFOTA 기능을 활성화/비활성화 할 수 있습니다.

\sphinxAtStartPar
\sphinxstylestrong{자동 업데이트 설정} : 새로운 펌웨어가 있는 경우 자동으로 업데이트를 진행합니다.


\subsection{AT 명령어}
\label{\detokenize{main_menu/modem:at}}
\noindent\sphinxincludegraphics[width=1.000\linewidth]{{atcommand-00}.png}

\sphinxAtStartPar
\sphinxstylestrong{AT 명령어 입력} : UI에서 AT 명령어를 입력하여, 모뎀과 통신할 수 있습니다.

\sphinxAtStartPar
\sphinxguilabel{명령어 전송} 실행시 하단에 응답 결과 에 결과가 표시 됩니다.


\subsection{모뎀 리셋}
\label{\detokenize{main_menu/modem:id7}}
\noindent\sphinxincludegraphics[width=0.800\linewidth]{{qreboot-00}.png}

\sphinxAtStartPar
Quectel 모뎀을 리셋합니다.

\begin{sphinxadmonition}{important}{중요:}
\sphinxAtStartPar
라우터 디바이스는 재시작 하지 않습니다.
\end{sphinxadmonition}

\sphinxstepscope


\section{네트워크}
\label{\detokenize{main_menu/network:id1}}\label{\detokenize{main_menu/network::doc}}

\bigskip\hrule\bigskip



\subsection{인터페이스}
\label{\detokenize{main_menu/network:id2}}
\noindent\sphinxincludegraphics[width=1.000\linewidth]{{interface-00}.png}

\sphinxAtStartPar
물리적 인터페이스의 현황을 보여줍니다. 간단한 사용량을 볼수 있으며, \sphinxguilabel{수정} 버튼을 눌러서 활성화/비활성화, 기본 게이트웨이 주소 변경 및 설정을 할수 있습니다.

\noindent\sphinxincludegraphics[width=0.700\linewidth]{{interface-01}.png}

\begin{sphinxadmonition}{caution}{조심:}
\sphinxAtStartPar
기본 게이트웨이 주소가 올바르지 못하면, UI에 접속이 불가할수 있으며, 공장 초기화(factory reset)을 해야 할수 있습니다.
\end{sphinxadmonition}


\subsection{장치}
\label{\detokenize{main_menu/network:id3}}
\noindent\sphinxincludegraphics[width=1.000\linewidth]{{interface-02}.png}

\sphinxAtStartPar
논리적인 인터페이스의 현황을 보여줍니다. \sphinxguilabel{설정} 버튼을 눌러서 MTU 값을 수정할 수 있습니다.

\noindent\sphinxincludegraphics[width=0.700\linewidth]{{interface-03}.png}


\subsubsection{DHCP}
\label{\detokenize{main_menu/network:dhcp}}
\sphinxAtStartPar
현재 접속되어 있거나 접속 되었던 DHCP 호스트 들의 정보를 표시 해줍니다. 설정을 통해서 기본 게이트웨이 주소를 변경할수 있으며, DHCP 범위도 변경이 가능합니다.


\subsection{할당}
\label{\detokenize{main_menu/network:id4}}
\noindent\sphinxincludegraphics[width=1.000\linewidth]{{dhcp-00}.png}

\sphinxAtStartPar
\sphinxguilabel{추가} 혹은 \sphinxguilabel{수정} 버튼을 누르면 아래 화면 처럼 DHCP 등록을 할 수 있는 화면이 표시 됩니다.

\noindent\sphinxincludegraphics[width=1.000\linewidth]{{dhcp-01}.png}


\subsection{설정}
\label{\detokenize{main_menu/network:id5}}
\sphinxAtStartPar
기본 게이트웨이 주소를 변경할수 있으며, DHCP 범위도 변경이 가능합니다.

\noindent\sphinxincludegraphics[width=1.000\linewidth]{{dhcp-02}.png}

\begin{sphinxadmonition}{caution}{조심:}
\sphinxAtStartPar
기본 게이트웨이 주소가 올바르지 못하면, UI에 접속이 불가할수 있으며, 공장 초기화(factory reset)을 해야 할수 있습니다.
\end{sphinxadmonition}


\subsubsection{방화벽}
\label{\detokenize{main_menu/network:id6}}
\begin{sphinxadmonition}{important}{중요:}
\sphinxAtStartPar
CGNAT(Carrier\sphinxhyphen{}Grade NAT) 구간에서의 사용은, 동작에 물리적인 제한이 있습니다. 예) inbound 처리
\end{sphinxadmonition}

\begin{sphinxadmonition}{caution}{조심:}
\sphinxAtStartPar
방화벽 설정이 올바르지 못하면, UI에 접속이 불가할 수 있으며, 공장 초기화(factory reset)을 해야 할수 있습니다.
\end{sphinxadmonition}


\subsection{일반 설정}
\label{\detokenize{main_menu/network:id7}}
\noindent\sphinxincludegraphics[width=1.000\linewidth]{{firewall-00}.png}

\sphinxAtStartPar
방화벽은 네트워크 인터페이스를 영역(Zone)으로 묶어 네트워크 트래픽 흐름을 제어합니다.


\subsection{포트 포워딩}
\label{\detokenize{main_menu/network:id8}}
\noindent\sphinxincludegraphics[width=1.000\linewidth]{{firewall-01}.png}

\sphinxAtStartPar
포트 포워딩을 통해 인터넷상의 외부 컴퓨터가 내부 LAN의 특정 컴퓨터나 서비스에 접속할 수 있습니다.


\subsection{트래픽 규칙}
\label{\detokenize{main_menu/network:id9}}
\noindent\sphinxincludegraphics[width=1.000\linewidth]{{firewall-02}.png}

\sphinxAtStartPar
트래픽 규칙은 서로 다른 영역(Zone) 사이를 오가는 패킷에 대한 정책을 정의합니다. 예를 들어 특정 호스트 간의 통신을 차단하거나 라우터의 WAN 포트를 개방할 수 있습니다.


\subsection{NAT 규칙}
\label{\detokenize{main_menu/network:nat}}
\noindent\sphinxincludegraphics[width=1.000\linewidth]{{firewall-03}.png}

\sphinxAtStartPar
NAT 규칙을 사용하여 외부로 나가거나 전달되는 트래픽의 출발지 IP를 세밀하게 제어할 수 있습니다.


\subsection{DMZ}
\label{\detokenize{main_menu/network:dmz}}
\noindent\sphinxincludegraphics[width=1.000\linewidth]{{firewall-04}.png}

\sphinxAtStartPar
모든 포트를 열고 단일 로컬 장치를 인터넷에 완전히 노출하는 기능 입니다.

\sphinxAtStartPar
CGNAT(Carrier\sphinxhyphen{}Grade NAT)구간에서의 사용은, 동작에 제한이 있으며, 만약 WAN ip주소가 해당 대역대로 확인되면 아래와 같은 경고창이 표시 됩니다.

\sphinxAtStartPar
\sphinxstylestrong{현재 LTE 연결이 사설 IPv4 주소를 사용하고 있습니다. DMZ를 활성화해도 인터넷에서 들어오는 연결이 불가능할 수 있습니다.}
\begin{enumerate}
\sphinxsetlistlabels{\arabic}{enumi}{enumii}{}{.}%
\item {} 
\sphinxAtStartPar
\sphinxstylestrong{DMZ 사용} :  만약 사용을 원하면 해당 내용을 체크 합니다.

\item {} 
\sphinxAtStartPar
\sphinxstylestrong{DMZ 호스트 IP} : 사용을 원하는 디바이스가 있을시에 해당 IP를 설정 합니다.

\end{enumerate}

\begin{sphinxadmonition}{caution}{조심:}
\sphinxAtStartPar
가능하면 DMZ 대신 표준 포트 포워딩을 사용하여 애플리케이션에 필요한 특정 포트만 여십시오
\end{sphinxadmonition}

\sphinxstepscope


\section{시스템}
\label{\detokenize{main_menu/system:id1}}\label{\detokenize{main_menu/system::doc}}

\bigskip\hrule\bigskip



\subsection{일반 설정}
\label{\detokenize{main_menu/system:id2}}
\noindent\sphinxincludegraphics[width=0.800\linewidth]{{system-00}.png}
\begin{enumerate}
\sphinxsetlistlabels{\arabic}{enumi}{enumii}{}{.}%
\item {} 
\sphinxAtStartPar
\sphinxstylestrong{로컬 시간} : 현재 디바이스의 시간을 확인할 수 있습니다.

\end{enumerate}

\sphinxAtStartPar
\sphinxguilabel{브라우저 시간대로 동기화}: 현재 사용중인 환경에 브라우저의 시간에 맞춥니다.

\sphinxAtStartPar
\sphinxguilabel{NTP 서버와 동기화} : 시간 동기화 메뉴에 설정된 NTP 서버와 동기화 합니다.
\begin{enumerate}
\sphinxsetlistlabels{\arabic}{enumi}{enumii}{}{.}%
\setcounter{enumi}{1}
\item {} 
\sphinxAtStartPar
\sphinxstylestrong{호스트 이름}: 라우터 디바이스의 호스트명을 정할수 있습니다.

\item {} 
\sphinxAtStartPar
\sphinxstylestrong{설명}: 디바이스에 대한 설명을 입력할 수 있습니다.(선택사항)

\item {} 
\sphinxAtStartPar
\sphinxstylestrong{메모}: 디바이스에 대한 메모를 입력할 수 있습니다.(선택사항)

\item {} 
\sphinxAtStartPar
\sphinxstylestrong{시간대}: 디바이스의 시간대를 설정할 수 있습니다.

\item {} 
\sphinxAtStartPar
\sphinxstylestrong{시간 형식}: 디바이스의 시간 형식을 설정할 수 있습니다. 12시간제와 24시간제를 지원합니다.

\end{enumerate}


\subsection{시간 동기화}
\label{\detokenize{main_menu/system:id3}}
\noindent\sphinxincludegraphics[width=0.800\linewidth]{{ntp-00}.png}

\sphinxAtStartPar
NTP에 관련한 설정을 할 수 있습니다.

\sphinxAtStartPar
NTP 서버 후보에 NTP 서버 주소를 입력하고, \sphinxguilabel{저장 및 적용} 버튼을 클릭하면 NTP 서버와 동기화가 수행됩니다.


\subsection{언어 및 스타일}
\label{\detokenize{main_menu/system:id4}}
\noindent\sphinxincludegraphics[width=0.800\linewidth]{{theme-00}.png}

\sphinxAtStartPar
언어와 스타일을 설정 합니다. 언어는 \sphinxstylestrong{한국어} 와 \sphinxstylestrong{영어} 를 지원하며, 스타일은 \sphinxstylestrong{NTmore} 기본 theme를 제공합니다.


\subsubsection{관리}
\label{\detokenize{main_menu/system:id5}}

\subsection{라우터 암호}
\label{\detokenize{main_menu/system:id6}}
\noindent\sphinxincludegraphics[width=0.800\linewidth]{{password-00}.png}

\sphinxAtStartPar
디바이스의 패스워드를 변경 합니다.


\subsection{HTTP(S) 접속}
\label{\detokenize{main_menu/system:http-s}}
\noindent\sphinxincludegraphics[width=0.800\linewidth]{{https-00}.png}

\sphinxAtStartPar
관리화면 접속시 HTTP 접속을 HTTPS 접속으로 리디이렉션 합니다.


\subsubsection{시스템 업데이트}
\label{\detokenize{main_menu/system:id7}}
\sphinxAtStartPar
디바이스에 대한 업데이트 설정을 확인하고, 새로운 펌웨어가 있는 경우 업데이트를 수행할 수 있습니다.


\subsection{정보}
\label{\detokenize{main_menu/system:id8}}
\noindent\sphinxincludegraphics[width=0.800\linewidth]{{fota-00}.png}
\begin{enumerate}
\sphinxsetlistlabels{\arabic}{enumi}{enumii}{}{.}%
\item {} 
\sphinxAtStartPar
\sphinxstylestrong{기기 모델}

\item {} 
\sphinxAtStartPar
\sphinxstylestrong{디바이스 일련번호}

\item {} 
\sphinxAtStartPar
\sphinxstylestrong{현재 펌웨어 버전} : 현재 펌웨어 버전을 확인할 수 있습니다. (revision/build date) 형식으로 구성 되어 있습니다.

\item {} 
\sphinxAtStartPar
\sphinxstylestrong{FOTA 연동상태} : 현재 FOTA 서버와의 연동 상태를 확인할 수 있습니다.

\end{enumerate}

\sphinxAtStartPar
\sphinxguilabel{재부팅 실행} 버튼을 클릭하면 FOTA 서버와의 연동 상태를 확인할 수 있습니다.

\sphinxAtStartPar
\sphinxguilabel{즉시 확인} 버튼을 클릭하면 FOTA 서버에서 새로운 펌웨어가 있는지 확인할 수 있습니다.’


\subsection{설정}
\label{\detokenize{main_menu/system:id9}}
\noindent\sphinxincludegraphics[width=0.800\linewidth]{{fota-01}.png}
\begin{enumerate}
\sphinxsetlistlabels{\arabic}{enumi}{enumii}{}{.}%
\item {} 
\sphinxAtStartPar
\sphinxstylestrong{FOTA 기능 활성화} : FOTA 펌웨어 업데이트 기능을 사용합니다. 만약 FOTA 기능을 사용하지 않으려면 체크를 해제하세요.

\item {} 
\sphinxAtStartPar
\sphinxstylestrong{자동 업데이트 설정} : FOTA 서버에서 새로운 펌웨어가 있는 경우 자동으로 업데이트를 수행합니다. 만약 자동 업데이트를 사용하지 않으려면 체크를 해제하세요.

\end{enumerate}


\subsubsection{Backup / Restore config}
\label{\detokenize{main_menu/system:backup-restore-config}}
\noindent\sphinxincludegraphics[width=1.000\linewidth]{{backup-restore-00}.png}


\subsection{작업}
\label{\detokenize{main_menu/system:id10}}\begin{enumerate}
\sphinxsetlistlabels{\arabic}{enumi}{enumii}{}{.}%
\item {} 
\sphinxAtStartPar
\sphinxstylestrong{백업}

\end{enumerate}
\begin{itemize}
\item {} 
\sphinxAtStartPar
현재 설정 파일에 대한 tar 아카이브 다운로드를 원한다면 \sphinxguilabel{아카이브 생성} 버튼을 클릭하세요.

\item {} 
\sphinxAtStartPar
백업된 파일은 브라우저를 통해 PC에 다운로드 됩니다.

\end{itemize}
\begin{enumerate}
\sphinxsetlistlabels{\arabic}{enumi}{enumii}{}{.}%
\setcounter{enumi}{1}
\item {} 
\sphinxAtStartPar
\sphinxstylestrong{복원}

\end{enumerate}
\begin{itemize}
\item {} \begin{description}
\sphinxlineitem{\sphinxstylestrong{기본값으로 초기화}}\begin{itemize}
\item {} 
\sphinxAtStartPar
설정 파일을 초기 상태로 되돌리려면 \sphinxcode{\sphinxupquote{초기화 수행}} 을 클릭하세요.

\item {} 
\sphinxAtStartPar
초기화 수행후에는 시스템이 재부팅 됩니다.

\end{itemize}

\end{description}

\item {} \begin{description}
\sphinxlineitem{\sphinxstylestrong{백업 복원}}\begin{itemize}
\item {} 
\sphinxAtStartPar
설정 파일을 복원하려면 이전에 생성한 백업 아카이브를 \sphinxguilabel{아카이브 업로드} 에 업로드 하세요.

\item {} 
\sphinxAtStartPar
재부팅후에는 시스템이 재부팅 됩니다.

\end{itemize}

\end{description}

\end{itemize}


\subsubsection{재부팅}
\label{\detokenize{main_menu/system:id11}}
\noindent\sphinxincludegraphics[width=1.000\linewidth]{{reboot-00}.png}

\sphinxAtStartPar
디바이스를 재시작 합니다. \sphinxguilabel{재부팅 실행} 버튼을 클릭하면 시스템이 재부팅 됩니다.

\sphinxstepscope


\section{도움말}
\label{\detokenize{main_menu/help:id1}}\label{\detokenize{main_menu/help::doc}}
\noindent\sphinxincludegraphics[width=1.000\linewidth]{{help-00}.png}

\sphinxAtStartPar
현재 디바이스에 대한 버전 정보를 표시 합니다.

\sphinxAtStartPar
문서에 대한 내용은 \sphinxurl{https://ntmore-ntlr800z-b1-doc.readthedocs.io} 에서도 확인할 수 있습니다.


\subsection{오픈소스 라이선스}
\label{\detokenize{main_menu/help:id2}}
\noindent\sphinxincludegraphics[width=1.000\linewidth]{{help-01}.png}

\sphinxAtStartPar
사용한 오픈소스 라이선스에 대한 정보를 표시 합니다.

\sphinxstepscope


\section{로그아웃}
\label{\detokenize{main_menu/logout:id1}}\label{\detokenize{main_menu/logout::doc}}

\bigskip\hrule\bigskip


\sphinxAtStartPar
\sphinxstylestrong{로그아웃} 메뉴를 누를경우 바로 로그아웃이 진행되며, 로그인 화면으로 이동합니다.



\renewcommand{\indexname}{색인}
\printindex
\end{document}